% ============================================================
%  part1/ch02-photograph.tex
%  Chapter 2: The Photograph
% ============================================================

\chapter{The Photograph}
\label{ch:photograph}

\epigraph{%
  All photographs are accurate.\\
  None of them is the truth.%
}{Richard Avedon}

\section{A more faithful record}

A photograph feels more faithful than a map.

A road map replaces a city with coloured lines.
A photograph at least looks like the place it depicts.
The proportions are approximately right.
The colours are approximately right.
Faces are recognisable.
Buildings are identifiable.

Yet the photograph destroys an enormous amount of information.

It destroys depth.
A photograph collapses three-dimensional space onto a
two-dimensional surface.
Objects at different distances appear at the same depth.

It destroys time.
A photograph captures a single moment from a continuous process.
The moment before and the moment after are lost.

It destroys temperature, smell, sound, motion, and the
texture of surfaces.

It destroys causal history.
The photograph of a building does not record who built it,
what stood there before, or what decisions led to its current form.

The photograph feels faithful because it preserves something that
maps often discard: the visual geometry of a scene.
But it is still a projection, and its kernel is vast.

\section{The projection and its kernel}

Let the world be modeled as a four-dimensional event field:
\[
  W \subset \mathbb{R}^3 \times \mathbb{R}.
\]
A photograph is a projection onto an image plane:
\[
  P : W \to \mathbb{R}^2.
\]
If $I = P(W)$ denotes the resulting image, then many different
worlds can produce the same image:
\[
  P(W_1) = P(W_2) = I, \qquad W_1 \neq W_2.
\]

\begin{example}[Non-injectivity of photography]
  Two scenes are photographically equivalent if they produce the
  same image from a given camera position.
  A flat painting of a room and an actual room, viewed from the
  same angle, can produce identical photographs.
  A large object far away and a small object nearby can subtend
  the same visual angle.
  A moving object and a stationary one can appear identical in a
  sufficiently short exposure.
\end{example}

The inverse problem
\[
  P^{-1}(I)
\]
is therefore not a single world but an entire space of compatible
worlds.

\begin{definition}[Admissible reconstruction class]
  Given an image $I \in \mathbb{R}^2$, the
  \emph{admissible reconstruction class} of $I$ is
  \[
    [I]_P = \{W \in \mathcal{W} : P(W) = I\},
  \]
  where $\mathcal{W}$ is the space of possible world-states.
  A photograph does not reveal a unique world.
  It defines an equivalence class of worlds consistent with the
  observed image.
\end{definition}

This is the first appearance of what will later become the central
technical concept of the book: the admissibility class.

A projection does not select a unique preimage.
It selects a set of preimages.
Reconstruction requires additional conditions to choose among them.

\section{Non-injectivity and the inverse problem}

The failure of injectivity is not incidental.
It is the defining feature of observation.

\begin{proposition}[Non-injectivity of observation]
  Let $\pi : X \to M$ be any observation map where $\dim X >
  \dim M$.
  Then $\pi$ is non-injective: there exist $x_1 \neq x_2$ in $X$
  with $\pi(x_1) = \pi(x_2)$.
\end{proposition}

\begin{proof}
  By dimension counting: the fibres $\pi^{-1}(m)$ for generic
  $m \in M$ are $({\dim X - \dim M})$-dimensional submanifolds of
  $X$.
  Since $\dim X > \dim M$, these fibres are nontrivial.
\end{proof}

Every observation maps from a higher-dimensional reality to a
lower-dimensional record.
Non-injectivity is therefore unavoidable whenever reality is
richer than observation, which is always.

\section{Reconstruction as the central problem}

Given a photograph $I$, an observer attempting to reconstruct the
world must find a $W \in [I]_P$ that is consistent with all
available information.

This is an inverse problem:
\[
  \text{Find } W \text{ such that } P(W) = I
  \text{ and } W \text{ is admissible.}
\]
Admissibility here means consistent with prior knowledge, physical
constraints, and other available observations.

\begin{example}[Depth reconstruction from stereo vision]
  The human visual system solves this inverse problem
  continuously.
  Two slightly different photographs --- one from each eye ---
  combined with prior knowledge about the scale of familiar objects
  allows the visual system to reconstruct approximate depth.
  The reconstruction is not guaranteed to be correct.
  It is guaranteed to be admissible given the available data.
\end{example}

The reconstruction error measures how far the recovered world
departs from the true one:
\[
  E(W) = d(W, \hat{W}),
\]
where $\hat{W} = G(P(W))$ is the reconstructed world and
$d$ is some appropriate metric on world-states.

This is the first appearance of the reconstruction error
$E(x) = d(x, G(F(x)))$ that will become central in
Chapter~\ref{ch:memory} and throughout Part~V\@.

\section{What the photograph reveals about admissibility}

The photograph chapter advances the argument in one important
direction beyond the map chapter.

The map chapter showed that every observation loses information.

The photograph chapter shows that the lost information defines
an equivalence class of possible originals.
The question is no longer merely: what is missing?
The question becomes: among all originals consistent with the
observation, which are admissible?

\principle{
  A photograph does not reveal reality.\\[4pt]
  It defines a class of realities consistent with the
  observation.\\[4pt]
  Reconstruction selects from this class using admissibility
  conditions.
}

This principle applies to every observation system encountered in
this book: memories, detectors, embeddings, horizons, and
cosmological surveys.

The word \emph{admissibility} will be made mathematically precise
in Part~V\@.
For now, it means: consistent with the observation and with
whatever prior constraints are available.

\section{A note on the photograph and time}

One dimension the photograph discards deserves special attention:
time.

A photograph freezes a process into a state.
The process itself --- the causal history that produced the state
--- is not recorded.
Yet causal history is precisely what determines which features of
the image are meaningful.

A crack in a wall means one thing if it appeared yesterday and
another if it has been there for a century.
A face looks different knowing it belongs to someone young rather
than old.

The photograph loses trajectory information.
It preserves only a single state.

This will become important in Chapter~\ref{ch:detector} and
throughout Part~II, where trajectory information --- not just
states --- turns out to be essential for understanding dynamical
systems.

The first hint of the RSVP framework's emphasis on paths and
histories rather than states is already visible in the simple
observation that a photograph destroys time.
